\documentclass{article}
\usepackage[T1]{fontenc}      
\usepackage[utf8]{inputenc}    
\usepackage{lmodern}           
\usepackage{amsmath, amssymb, amsthm}  % Maths + proof
\usepackage{mathtools}         
\usepackage{geometry}
\theoremstyle{remark} 
\newtheorem*{remark}{Remarque}          
\geometry{margin=2.5cm}

\title{Fonctions mesurables et tribus}
\author{David Thébault}
\date{\today}

% Définition des environnements théorèmes
\theoremstyle{definition}
\newtheorem{definition}{Définition}[section]
\newtheorem{example}{Exemple}[section]

\theoremstyle{plain}
\newtheorem{theorem}{Théorème}[section]
\newtheorem{proposition}{Proposition}[section]

\begin{document}
\maketitle

\section{Tribus et tribu borélienne}

\begin{definition}[Tribu]
Soit $X$ un ensemble. Une \emph{tribu} $\mathcal A$ sur $X$ est une famille de parties
de $X$ telle que :
\begin{enumerate}
    \item $X \in \mathcal A$,
    \item $A \in \mathcal A \implies A^c \in \mathcal A$,
    \item $(A_n)_{n\ge1} \subset \mathcal A \implies \bigcup_{n\ge1} A_n \in \mathcal A$.
\end{enumerate}
\end{definition}

\begin{definition}[Tribu borélienne]
La \emph{tribu borélienne} de $\mathbb R$, notée $\mathcal B(\mathbb R)$, est la plus petite
tribu contenant tous les ouverts de $\mathbb R$.
\end{definition}

\begin{proposition}
La tribu borélienne de $\mathbb{R}$ est engendrée par les intervalles de la forme $(-\infty,a]$ :
\[
\mathcal{B}(\mathbb{R}) = \sigma\big((-\infty,a] \; ; \; a \in \mathbb{R}\big).
\]
\end{proposition}

\begin{proof}
Posons $\mathcal{G} = \{(-\infty,a]\; ;\; a \in \mathbb{R}\}$.  
Nous montrons les deux inclusions.

\textbf{1. $\sigma(\mathcal G) \subset \mathcal B(\mathbb R)$ :}  
Pour tout $a \in \mathbb R$, $(-\infty,a] = \bigcap_{n=1}^{\infty} (-\infty, a+1/n)$,  
et chaque intervalle $(-\infty,a+1/n)$ est ouvert. Comme $\mathcal B(\mathbb R)$ est une tribu, $(-\infty,a] \in \mathcal B(\mathbb R)$. Donc $\mathcal G \subset \mathcal B(\mathbb R)$ et $\sigma(\mathcal G) \subset \mathcal B(\mathbb R)$.

\textbf{2. $\mathcal B(\mathbb R) \subset \sigma(\mathcal G)$ :}  
Pour $a<b$, $(a,b) = (-\infty,b) \cap (-\infty,a]^c$, et
$(-\infty,b) = \bigcup_{n=1}^{\infty} (-\infty,b-1/n] \in \sigma(\mathcal G)$.
Comme $(-\infty,a]^c \in \sigma(\mathcal G)$ et $\sigma(\mathcal G)$ est stable par intersection, $(a,b) \in \sigma(\mathcal G)$.  
Tout ouvert étant une réunion d'intervalles, tous les ouverts appartiennent à $\sigma(\mathcal G)$, donc $\mathcal B(\mathbb R) \subset \sigma(\mathcal G)$.
\end{proof}

\section{Fonctions mesurables}

\begin{definition}[Image réciproque]
Soit $f : X \to Y$ et $A \subset Y$.  
L'image réciproque de $A$ par $f$ est
\[
f^{-1}(A) = \{x \in X \mid f(x) \in A\}.
\]
Cette notion est définie pour toute fonction, indépendamment de toute injectivité ou inversibilité.
\end{definition}

\begin{proposition}[Propriétés de l'image réciproque]
Pour tous ensembles $A, (A_n)_{n\ge1} \subset Y$ :
\[
f^{-1}(A^c) = (f^{-1}(A))^c,
\qquad
f^{-1}\Big(\bigcup_{n\ge1} A_n\Big) = \bigcup_{n\ge1} f^{-1}(A_n).
\]
\end{proposition}

\begin{proof}
Pour tout $x \in X$ :
\[
x \in f^{-1}(A^c) \iff f(x) \in A^c \iff f(x) \notin A \iff x \notin f^{-1}(A) \iff x \in (f^{-1}(A))^c.
\]

De même,
\[
x \in f^{-1}\Big(\bigcup_n A_n\Big) \iff f(x) \in \bigcup_n A_n \iff \exists n, f(x) \in A_n \iff x \in \bigcup_n f^{-1}(A_n).
\]
\end{proof}

\begin{definition}[Fonction mesurable]
Soient $(X,\mathcal A)$ et $(Y,\mathcal B)$ deux espaces mesurables.  
Une fonction $f : X \to Y$ est dite \emph{mesurable} si
\[
\forall B \in \mathcal B, \quad f^{-1}(B) \in \mathcal A.
\]
\end{definition}

\begin{theorem}[Fonctions réelles mesurables]
Soit $f : \mathbb R \to \mathbb R$, où $\mathbb R$ est muni de la tribu borélienne.  
Alors $f$ est mesurable si et seulement si
\[
\forall a \in \mathbb R, \quad \{x \in \mathbb R \mid f(x) \le a \} \in \mathcal B(\mathbb R).
\]
\end{theorem}

\begin{proof}
Les ensembles $(-\infty,a]$ engendrent la tribu borélienne et l'image réciproque respecte
les opérations de tribu (union et complémentaire). Il suffit donc de vérifier la condition pour cette famille génératrice.
\end{proof}

\begin{remark}
Une fonction mesurable garantit que pour tout ensemble borélien $B \subset \mathbb R$,  
l'ensemble $\{x \mid f(x) \in B\}$ est mesurable.  

En probabilité, une variable aléatoire réelle est précisément une fonction
\[
X : (\Omega, \mathcal F) \to (\mathbb R, \mathcal B(\mathbb R)),
\]  
ce qui garantit que les événements $\{X \le a\}$ sont bien définis pour tout $a \in \mathbb R$.
\end{remark}

\end{document}